% ------------------------------------------------------------------------------
% Chapter 2
% ------------------------------------------------------------------------------
\chapter{Literature Review} % enter the name of the chapter here
\label{cha:ch2_literature} % enter the chapter label here (for cross referencing)

% ------------------------------------------------------------------------------
% Introduction
% ------------------------------------------------------------------------------
\section{Introduction}
\label{sec:Introduction}

Generative AI is a structural inflection point in knowledge-intensive work. It is different from other forms of AI and automation because the former primarily focuses on complex cognitive work through the creation of outputs that are sensitive to context and require interpretation, refinement, and integration. In the field of software engineering, AI coding assistants in the form of generative code completion tools are embedded into the coding environment and allow for the generation, debugging, documentation, and refactoring of code.

Empirical studies of the impact of generative AI on human productivity and output quality revealed that, under conditions of structured task design, a significant increase in human output is achieved. In a large-scale experimental design, the availability of generative AI tools resulted in a reduction in task completion time and an increase in output quality, as determined by the researcher. The results achieved were statistically significant for the total sample, while the results were more pronounced for low-performing individuals \autocite{Noy2023}. The results support the idea that generative AI tools result in micro-level productivity effects.

However, the productivity gains at the task level do not necessarily translate into organizational return on investment. Research on AI investment at the firm level indicates that the performance effects are heterogeneous and conditional. AI investment is linked with growth through innovation and product development channels; however, the level and stability of the returns depend on the capabilities of the organization \autocite{Babina2024}. The market reaction to AI investment announcements is mixed, indicating the level of awareness among investors about the risk and capabilities \autocite{Lui2021}.

These studies all circle back to the same basic idea: AI capability does not produce value; rather, value is produced when the technological tools are embedded in the organizational systems. More recent studies have suggested another component: a digital organizational culture can influence the innovation outcome of AI capability \autocite{An2024}. Leadership can also play a part in influencing the development of trust and embedding AI usage \autocite{Schneider2023}.

Despite the recent increase in research on the performance effects of AI, little research has appeared that combines the effects of organizational culture and leadership into a structured ROI model for AI coding assistants in software development teams. This chapter synthesizes global empirical research on AI productivity and firm performance, critically evaluates the impact of organizational conditions, and focuses the research question more specifically on the South African market, where empirical ROI research has yet to be conducted.

% -----------------------------------------------------------------------------------
% Generative AI and Productivity: From Task Performance to Organisational Outcomes
% -----------------------------------------------------------------------------------
\section{Generative AI and Productivity: From Task Performance to Organisational Outcomes}
%\label{sec:Generative AI and Productivity: From Task Performance to Organisational Outcomes}

The best evidence for how generative AI changes work comes from tightly controlled 
experiments. In the case of \textcite{Noy2023}, the evidence provides strong causal proof of how gen AI increases the rate of task completion and enhances the quality of professional writing. From this, we can see precisely how the presence of gen AI changes the gap in performance. The benefits of using gen AI are not equally distributed, with the less skilled benefiting the most. This indicates the ability of gen AI to level the playing field for teams.

This uneven benefit to different individuals can, therefore, have implications for software 
development teams. If the benefits to gen AI for coding tasks are greater for the less skilled, the composition of the team might change. However, the benefits to the team must not only be seen at the task level. They must be sustained, integrated, and monetized to increase the ROI.

This distinction finds corroboration in firm-level evidence. \textcite{Babina2024} show a 
positive association between the adoption of AI and firm growth and innovation, with the 
Association is stronger for firms with complementary capabilities. The channel through which this association seems to occur is mainly through product innovation rather than labour substitution effects, implying the importance of AI in value creation through product development rather than cost reduction.

At the same time, there is evidence of mixed reactions in the financial markets to AI 
investment announcements. \textcite{Lui2021} show that some firms have positive abnormal returns around the announcement of AI investments, whereas others have neutral or negative returns. This appears to reflect investors' ability to distinguish between firms with the ability to successfully implement AI and those with a risk of implementation failure.

In the context of software engineering, empirical evidence points to the fact that there is a great deal of usage of AI coding assistants. However, this usage is done in an inconsistent manner. In a comprehensive empirical study done to examine the experiences of software developers, it was found that there was usage of AI coding assistants for various activities in the software development process. However, issues of trust, boundaries, and governance ambiguity have been found to hinder the effective impact of this usage \autocite{Sergeyuk2025}. This points to the fact that the frequency of usage alone may not guarantee the achievement of benefits.

According to the literature, the return on investment (ROI) model of AI coding assistants comprises a three-tier structure:

\subsection {Effects on productivity at the task level.}
\label{sec:Effects on productivity at the task level}
Task-level productivity effects relate to the observed improvements that accrue when individual programmers utilize AI-based coding assistants for the accomplishment of specific programming tasks. At the task level, AI-based coding assistants are used for activities such as code writing, documentation, debugging, and suggestion of solutions to programming problems. Such activities have the potential to hasten the accomplishment of specific tasks through the provision of real-time support during the entire process. According to experimental research, generative AI-based tools have the capability to hasten the execution of knowledge-based tasks by a large factor without compromising the quality of the output \autocite{Noy2023}. This implies that, at the organizational level, the ability of the development team to hasten the accomplishment of specific coding tasks has the capability to hasten the accomplishment of more complex development tasks.

\subsection {Effects on integration and coordination at the workflow level.}
\label{sec:Effects on integration and coordination at the workflow level}
Workflow-level integration and coordination effects refer to the impact of AI coding assistants in the collaboration and coordination of teams in the development of software applications. Software development involves the implementation of workflows that include code reviews, testing procedures, version control, and collaboration between different developers. The incorporation of AI coding assistants in the development process can have positive effects on the coordination and collaboration of teams in the development process. This means that the teams can be able to achieve the required objectives in the most efficient manner possible. The benefits of using AI coding assistants have been examined in the experiences of developers using the system, and the findings suggest that the benefits of the system depend on the extent of its integration in the development process \autocite{Sergeyuk2025}.

\subsection {Effects on the firm's business finance and innovation at the firm level.}
\label{sec:Effects on the firm's business finance and innovation at the firm level}
Firm-level financial and innovation results refer to the overall organizational benefits that are realized when the benefits of increased productivity and workflow improvements are realized at the financial level. At the firm level, the organization seeks to determine whether the implementation of AI technologies results in benefits such as improved operational efficiency, lower development costs, faster time-to-market, or new forms of digital products/services, among others. Empirical studies conducted to determine the impact of AI technology adoption by organizations revealed that artificial intelligence technologies could drive organizational growth through innovation and improved productivity \autocite{Babina2024}. Therefore, it is critical to understand the role of organizational culture and leadership in the successful implementation of AI coding assistants, given the fact that the achievement of a financial return is dependent on the organization’s ability to adopt the new technology.

The majority of the literature is found to have been limited to the first tier of the ROI model. However, the second and third tiers of the model have been found to be underdeveloped in the context of AI coding assistants. This underdevelopment of the second and third tiers of the model points to the fact that the effective usage and integration of AI coding assistants have been found to be mediating factors in the ROI model.	


% -----------------------------------------------------------------------------------
% Measuring ROI in AI Contexts: Financial, Operational, and Strategic Dimensions
% -----------------------------------------------------------------------------------
\section{Measuring ROI in AI Contexts: Financial, Operational, and Strategic Dimensions}
\label{sec:Measuring ROI in AI Contexts: Financial, Operational, and Strategic Dimensions}
Return on investment (ROI) is not absolute in nature but rather relative. ROI is generally described as net benefits earned divided by total costs incurred over a particular period. For AI initiatives, costs generally include subscription fees, integration costs, training costs, governance costs, different models used to mitigate risks, and workflow costs. 

The potential benefits that could be earned from implementing AI are as follows: 
\subsection {Time saved translates to reduced labour costs.}
\label{sec:Time saved translates to reduced labour costs}
A major advantage of AI-based coding assistants is the potential for these tools to decrease the time it takes for a programmer to do their job. If the time it takes to do development tasks is decreased, it implies that the labour costs associated with the development process will decrease. According to experimental research conducted on generative AI, the availability of AI-based assistance has a significant impact on the time it takes to complete a given knowledge-intensive task without compromising the quality of the output \autocite{Noy2023}. Similar results have been obtained by other researchers in the field of software engineering, where it was found that the availability of AI-based assistance for the development process can accelerate the development process \autocite{Peng2023}.

\subsection {Increased productivity.}
\label{sec:Increased productivity}
The other possible benefit associated with the use of AI coding assistants is related to the potential for enhancing development throughput. Throughput refers to the volume of work that a development team can process within a particular period. The potential for enhancing development throughput means that an organization can process a larger volume of software features, bug fixes, or updates within a single development cycle without a corresponding increase in the number of developers. Studies on the use of AI-based work tools indicate that generative AI tools have the potential to increase output levels substantially within a particular period, a factor that allows workers to process more deliverables within a particular period \autocite{Noy2023}.

\subsection {Defects reduced.}
\label{sec:Defects reduced}
The coding assistants that AI offers may be beneficial in improving the quality of software codes. Defect rates refer to the number of defects or errors that are associated with a given software. It is suggested that when a developer is offered coding assistance and intelligent coding suggestions, he or she may be able to avoid certain coding defects. Research studies have shown that AI-assisted software development is beneficial to developers since it offers coding accuracy and assistance in solving coding challenges. AI-assisted software development is beneficial since it offers coding accuracy and assistance in solving coding challenges \autocite{Peng2023}.

\subsection {Reduction in rework costs.}
\label{sec:Reduction in rework costs}
Rework occurs when software development teams are forced to dedicate additional resources to correcting errors, re-coding, and correcting defects that are detected after the development process. It is one of the key factors that contribute to costs in software development, considering its effect on development time and its ability to cause delays in software development. By improving coding precision and offering real-time support to software developers, there is an opportunity for coding assistants to reduce the need for rework. The timely detection of errors in the development process allows software development teams to avoid the significant costs that are likely to result from addressing errors in subsequent software development phases.

\subsection {Faster time to market.}
\label{sec:Faster time to market}
Another way in which the use of AI coding assistants may help reduce the time it takes to deliver new software products or features is through the reduction of the time to market. Time to market refers to the length of time it takes to develop a given software product or feature from the start of the project to the release of the feature to the market. The faster the time to market, the better the organization can respond to market opportunities. The empirical literature has proven that the use of AI in the business environment leads to faster innovation processes in the business environment \autocite{Babina2024}.

\subsection {Increased revenues via innovation.}
\label{sec:Increased revenues via innovation}
Finally, the adoption of AI has the potential to lead to organizational growth through innovation. This is because when an organization can build new features more easily or try new ideas more easily, it has the potential to create more revenue opportunities through digital products or services. The adoption of AI has the potential to lead to innovation-driven growth. The study on the adoption of AI by firms revealed that firms that adopted AI had a positive relationship between product innovation and financial performance \autocite{Babina2024}. This means that an AI coding assistant has the potential to lead to not only increased operational efficiency but also increased competitiveness of an organization.

Empirical research on AI adoption in organizations suggests that AI capability is a factor in innovation-driven growth \autocite{Babina2024}. For quantifying the potential benefits that could be earned from the adoption of AI in team-level software development, a more refined measurement of AI benefits is required. 

In the software development domain, cycle time, deployment frequency, defect density, and mean time to recovery act as a framework to assess the overall performance in the software development domain. These could be easily translated into monetary values using the cost of delay, defect costs, and resource allocation. 

The research on AI has generally used perception-based productivity metrics. Even in experiments, researchers have focused on short-term results rather than long-term monetary results \autocite{Noy2023}. 


\subsection {The ideal model for calculating ROI on AI should include: }
\label{sec:The ideal model for calculating ROI on AI should include}
\begin{itemize}
	\item \textbf{Performance Metrics} 
\end{itemize}
The objective engineering performance metrics are the measurable parameters used to evaluate the performance and productivity of software development teams. The importance of the performance metrics is undeniable since it provides empirical evidence of the performance of the AI coding assistants. The performance metrics used include the development cycle time, the deployment frequency, the defect density, and the mean time to recovery. The performance metrics are widely used to evaluate the performance of the software development process since it encompasses the productivity and quality aspects of the process. The performance of the software development process is essential since it provides evidence of the effectiveness of the AI coding assistants. For instance, the use of generative AI tools has been observed to improve the performance of the task completion speed and the quality of the outcome, which may be reflected in the performance of the software development process \autocite{Noy2023}. The importance of the performance metrics is undeniable since it provides empirical evidence of the performance of the AI coding assistants.

\begin{itemize}
	\item \textbf{Cost Calculations} 
\end{itemize}
Full cost accounting is a comprehensive measure of all costs associated with the adoption and maintenance of artificial intelligence technology. In terms of AI coding assistants, organizations must consider all costs beyond the costs associated with the software tool. These costs may include costs associated with training developers, tool integration, updating policies and procedures, and maintaining the environment that is necessary to support AI usage. There are also costs associated with oversight that organizations must consider when using AI technology. Research studies have shown that organizations tend to underestimate all costs associated with technology adoption when studies are based solely on direct technology costs \autocite{Babina2024}. To properly assess ROI, all costs associated with AI adoption must be considered. These costs include direct and indirect costs.

\begin{itemize}
	\item \textbf{Quality of organizational adoption} 
\end{itemize}
Organisational adoption quality relates to the extent and quality of the incorporation of AI tools in organisational routines and practices. It should be noted that the mere incorporation of AI coding assistants in the development environment does not guarantee the quality and extent of the usage of such tools in organisational settings. Organisational adoption quality of AI tools can be regarded as high if developers effectively incorporate AI tools in coding activities, testing activities, and collaborative activities such as code reviews. Empirical studies of the organisational adoption of AI tools suggest that the organisational benefits of AI technologies depend significantly on the extent and quality of the incorporation of these tools in organisational routines and activities and the extent of organisational learning \autocite{An2024}. Therefore, organisational adoption quality can be regarded as a critical factor in determining the organisational value of AI coding assistants.

\begin{itemize}
	\item \textbf{Governance Costs} 
\end{itemize}
Governance costs refer to the costs associated with managing and regulating the usage of artificial intelligence (AI) technologies in an organization. The usage of AI coding assistants poses various governance challenges to an organization in terms of data privacy, intellectual property rights, software quality assurance, and compliance with organizational policies. According to various scholarly research on the governance of AI, effective governance mechanisms are necessary to ensure that the usage of AI systems is within acceptable risk parameters and complies with organizational standards \autocite{Schneider2023}. However, implementing governance mechanisms requires time and administrative costs, which are additional costs to consider when evaluating the return on investment in AI.

\begin{itemize}
	\item \textbf{Risks Involved} 
\end{itemize}
The risk exposure adjustments relate to how potential risks associated with the adoption of AI are considered when evaluating financial returns. There is a potential risk of incorporating coding assistants that may lead to software defects, security breaches, or copyright issues. The adoption of AI is likely to increase software defects, security breaches, or copyright issues. This implies that an organisation is likely to incur additional costs when risks associated with AI adoption are realised. According to an empirical study on organisational adoption of AI, technological advantages need to be balanced with risk management and governance to ensure sustainable value creation \autocite{Schneider2023}. Adjusting the ROI formula to incorporate potential risks associated with AI adoption is essential to avoid overestimating potential benefits.

The lack of an integrated model in AI-based coding assistant research is a knowledge gap that this research aims to bridge.

\section {Organizational Culture as a Mediator of AI Value Realisation.}
\label{sec:Organizational Culture as a Mediator of AI Value Realisation}
Thus, organizational culture plays a significant role in the interpretation, trust, and institutionalization of AI tools. Digital organizational culture has been proven to enhance the link between AI capabilities and innovation performance \autocite{An2024}. In other words, the moderating effect shows that organizational culture plays a significant role in the payoff of technological capabilities.

Trust formation is a significant factor for the adoption of AI tools. \textcite{Daly2025} proved that the formation of trust with AI tools is a function of organizational experience and the presence of support systems and governance systems. Excessive trust is a threat to the adoption of AI tools, while a lack of trust is a barrier to the adoption of AI tools.

Psychological safety is a significant factor for the adoption of AI tools in the context of software engineering teams. In other words, the presence or absence of psychological safety is a significant factor for the transformation of AI suggestions into quality code.

While organizational culture is a significant factor for the adoption of AI tools, it is not quantified in the context of the AI ROI model. Most studies on the impact of organizational culture on AI adoption qualitatively address the impact of organizational culture on AI adoption. Thus, the present study proposes a quantified approach to organizational culture as a mediating factor for the impact of AI coding assistants on the ROI.

\section {Leadership and AI Governance.}
\label{sec:Leadership and AI Governance}
The leadership function assumes critical importance in ensuring that AI adoption is in sync with business goals and risk management processes. The body of research on AI governance has identified three interrelated but distinct layers of AI governance data governance, model governance, and organizational governance \autocite{Schneider2023}. However, leadership ensures that AI governance structures are embedded within the organization rather than being adopted in an informal and ad hoc manner.

The leadership function assumes importance in defining return on investment (ROI) in AI in the following areas:

\subsection {Executive sponsorship.}
\label{sec:Executive sponsorship}
Leadership has a pivotal role to play in ensuring the achievement of value from artificial intelligence technology. Research studies on the digital transformation have all proven the importance of leadership in the achievement of the objectives of technology adoption. Among the most important leadership constructs in the achievement of technology adoption objectives is executive sponsorship, which refers to the support role played by the organizational leadership in the achievement of the objectives of technology adoption. When the organizational leadership provides support to the objectives of technology adoption, it becomes easier to integrate emerging technology, such as the use of artificial intelligence coding assistants, into the organizational strategy instead of it being used as a standalone productivity tool by individual developers \autocite{Schneider2023}. Further research studies on artificial intelligence technology have proven the importance of the organizational leadership in the achievement of organizational preparedness to adopt digital innovations and the ability of employees to integrate artificial intelligence technology into their work practices \autocite{An2024}.

\subsection {Resource allocation.}
\label{sec:Resource allocation}
Another important role of a leader involves resource allocation, which includes the provisioning of the financial, technical, and human resources needed to support the effective integration of the various AI systems into the organizational processes. The effective adoption of AI systems in the organizational processes is more than just the acquisition of the various AI systems. For instance, empirical research on the economic impact of the adoption of AI systems has shown that the ability of the organization to invest in the acquisition of various capabilities would have a positive impact on the ability of the organization to realize productivity improvements from the various AI systems \autocite{Babina2024}. In the case of software engineering, the allocation of resources may limit the ability of the developer to experiment with the various coding assistants, thus limiting the ability to realize the various learning processes.

\subsection {Metrics alignment.}
\label{sec:Metrics alignment}
Leadership has a vital role to play in the achievement of metrics alignment, which refers to the alignment of organizational performance metrics and the anticipated outcomes of implementing AI systems. The coding assistants of AI systems influence various aspects of software development performance, including development speed, code quality, collaboration, and defect reduction. Studies on the productivity outcomes of generative AI systems suggest that, while generative AI systems are instrumental in expediting task completion, the performance outcomes are contingent on the evaluation and integration of the output with the organizational process \autocite{Noy2023}. In cases where speed is the main objective and not adequately balanced with quality outcomes, the adoption of generative AI systems may inadvertently result in an escalation of technical debt and defects. However, metrics that balance speed with quality outcomes allow an organization to leverage the performance outcomes of generative AI systems to achieve operational performance.

\subsection {Policy enforcement.}
\label{sec:Policy enforcement}
Another key responsibility of leaders in this case is policy enforcement, which is described as the development and implementation of governance systems that guide the usage of AI systems in the workflow of an organisation. The coding assistants in AI systems pose new challenges in terms of governance, including intellectual property rights and verification of the source of the code, and managing risks to security. The research on AI governance shows that when an organisation has a system of governance in place, it is in a better position to manage technological risks and sustain performance improvements with the adoption of AI systems \autocite{Schneider2023}. This reduces uncertainty among developers and provides guidance on the usage of AI systems.

\subsection {Ethical considerations.}
\label{sec:Ethical considerations}
Finally, the ethical considerations associated with the use of artificial intelligence need to be addressed by the leaders. In the case of generative artificial intelligence, ethical considerations arise regarding transparency, ownership of intellectual properties, data privacy, and bias. Ethical governance practices ensure the effective implementation of artificial intelligence technology in a manner that is consistent with the regulatory framework and ethical codes of practice. Ethical artificial intelligence practices help to build trust relationships with employees, which is a critical factor in the long-term institutionalization of artificial intelligence technology. Empirical research studies on the organizational responses to artificial intelligence technology have established the importance of ethical governance mechanisms in building trust relationships between employees and artificial intelligence technology, which is essential to sustain the integration of artificial intelligence technology into the organization \autocite{Daly2025}

If leadership were to be absent in AI adoption, then AI coding assistants would be used in an optional capacity rather than being used to drive business performance. In addition to this, it would also help in reducing risks in intellectual property and code authenticity.

However, research studies that have investigated leadership in AI adoption and its direct linkage with financial ROI in AI coding assistants are scarce. The body of research has primarily been cantered on AI governance principles but has seldom been able to quantify financial ROI.

This has led to the need to incorporate leadership as a key mediating variable in AI ROI modelling.

\section {South African Context and Emerging Market Considerations.}
\label{sec:South African Context and Emerging Market Considerations}
The process of AI adoption in EMs takes place in an institutional environment where capability heterogeneity and resource scarcity exist. In South Africa, for instance, the organizational stakeholders of the country's software engineering industry acknowledge the competitive benefits and labour adjustments of AI adoption \autocite{Poisat2024}.

South Africa's software engineering industry comprises large organizations with intricate systems and a strong demand for digital transformation. However, the governance level and the allocation of workforce competencies differ among organizations. This highlights the importance of leadership and organizational culture in the achievement of AI benefits. In addition, in resource-scarce environments, justifying the return on investment (ROI) of AI implementation becomes essential. In this case, misalignment in the organization increases the risk of implementation.

Despite the surge in interest in the adoption of AI in South African software engineering teams, there still exists a scarcity of peer-reviewed literature on the ROI of using AI coding assistants. This highlights the importance of this paper.

\section {Synthesis and Research Gap.}
\label{sec:Synthesis and Research Gap}
However, current studies indicate that generative AI models boost task-level productivity \autocite{Noy2023}. Moreover, AI investments are related to firm expansion and innovation outcomes \autocite{Babina2024}. On the contrary, market reactions to AI adoption exhibit heterogeneity, reflecting risks associated with AI adoption, like AI integration \autocite{Lui2021}. Furthermore, AI model performance is conditional upon organizational culture in the digital space \autocite{An2024}. The role of AI governance is emphasized, requiring alignment between leadership and strategic goals \autocite{Schneider2023}. In addition, contextual limitations exist in the South African environment \autocite{Poisat2024}.

Nevertheless, some critical knowledge gaps remain:
\begin{itemize}
	\item \textbf{Insufficient return on investment modelling for AI-based coding assistants} 
\end{itemize}
The current literature on the subject suggests that AI coding assistants can improve the productivity of developers. However, there is a lack of literature on the financial return on investment of AI coding assistants. Most of the literature focuses on the productivity of developers at the task level or the perceived usefulness of AI tools. For example, experimental studies have found that generative AI tools significantly improve the efficiency of task performance and the quality of outputs in a knowledge work setting \autocite{Noy2023}. Other studies on AI programming tools have found that developers can perform coding tasks quickly with the aid of AI tools \autocite{Peng2023}. However, there is a lack of literature on the financial return on investment of AI coding assistants.

\begin{itemize}
	\item \textbf{Inadequate inclusion of culture and leadership within the financial model} 
\end{itemize}
Organisational culture and leadership have been recognised as key factors in the achievement of successful technology implementation outcomes. The empirical findings have suggested that the organisational culture for digitalisation can strengthen the impact of the capabilities of artificial intelligence on innovation and performance \autocite{An2024}. Leadership has been observed as the key factor in the development of the governance structure for the implementation of responsible AI in the organisational environment \autocite{Schneider2023}. However, the importance of organisational culture and leadership has been recognised in the literature while studying these factors in isolation for the achievement of technological performance outcomes. Only limited empirical research has been conducted on the development of the organisational culture and leadership factors within the financial framework for the assessment of the return on investment in AI technologies.

\begin{itemize}
	\item \textbf{Insufficient measurement of the operation at the team level} 
\end{itemize}
The available literature on AI-driven productivity tends to focus more on individual performance measures and not team-level measures to ensure operational effectiveness. However, it is to be noted that software development is a team-based activity in which developers interact with each other through various processes. Studies that have been conducted to understand the role of AI coding assistants have shown that these tools are used at various stages of the software development lifecycle, though there is a significant variation in their usage \autocite{Sergeyuk2025}. It is not clear how these tools influence the productivity and quality of software engineering teams without analysing team-level measures.

\begin{itemize}
	\item \textbf{Lack of empirical research from the continent} 
\end{itemize}
Most of the research on AI adoption and its organizational outcomes is based on empirical research undertaken in North America, Europe, and parts of Asia. Therefore, there is a scarcity of empirical research on how these AI technologies impact organizational outcomes in Africa. Research on the impact of artificial intelligence on the workforce in South Africa has seen a surge of interest in artificial intelligence adoption while underscoring the need to conduct more research on its impact on organizational outcomes within the local environment \autocite{Poisat2024}. The scarcity of research on AI coding assistants on software engineering teams within South Africa impedes our understanding of how contextual factors impact the realization of the ROI on AI.

This research intends to fill the knowledge gap by developing and empirically testing an integrated model that connects culture, leadership, effective adoption, operation, and measurable return on investment for software development teams in South Africa.